%% Generated by Sphinx.
\def\sphinxdocclass{report}
\documentclass[letterpaper,10pt,spanish]{sphinxmanual}
\ifdefined\pdfpxdimen
   \let\sphinxpxdimen\pdfpxdimen\else\newdimen\sphinxpxdimen
\fi \sphinxpxdimen=.75bp\relax
\ifdefined\pdfimageresolution
    \pdfimageresolution= \numexpr \dimexpr1in\relax/\sphinxpxdimen\relax
\fi
%% let collapsible pdf bookmarks panel have high depth per default
\PassOptionsToPackage{bookmarksdepth=5}{hyperref}

\PassOptionsToPackage{booktabs}{sphinx}
\PassOptionsToPackage{colorrows}{sphinx}

\PassOptionsToPackage{warn}{textcomp}
\usepackage[utf8]{inputenc}
\ifdefined\DeclareUnicodeCharacter
% support both utf8 and utf8x syntaxes
  \ifdefined\DeclareUnicodeCharacterAsOptional
    \def\sphinxDUC#1{\DeclareUnicodeCharacter{"#1}}
  \else
    \let\sphinxDUC\DeclareUnicodeCharacter
  \fi
  \sphinxDUC{00A0}{\nobreakspace}
  \sphinxDUC{2500}{\sphinxunichar{2500}}
  \sphinxDUC{2502}{\sphinxunichar{2502}}
  \sphinxDUC{2514}{\sphinxunichar{2514}}
  \sphinxDUC{251C}{\sphinxunichar{251C}}
  \sphinxDUC{2572}{\textbackslash}
\fi
\usepackage{cmap}
\usepackage[T1]{fontenc}
\usepackage{amsmath,amssymb,amstext}
\usepackage{babel}



\usepackage{tgtermes}
\usepackage{tgheros}
\renewcommand{\ttdefault}{txtt}



\usepackage[Sonny]{fncychap}
\ChNameVar{\Large\normalfont\sffamily}
\ChTitleVar{\Large\normalfont\sffamily}
\usepackage{sphinx}

\fvset{fontsize=auto}
\usepackage{geometry}


% Include hyperref last.
\usepackage{hyperref}
% Fix anchor placement for figures with captions.
\usepackage{hypcap}% it must be loaded after hyperref.
% Set up styles of URL: it should be placed after hyperref.
\urlstyle{same}

\addto\captionsspanish{\renewcommand{\contentsname}{Contents:}}

\usepackage{sphinxmessages}
\setcounter{tocdepth}{1}



\title{Tienda\_online}
\date{10 de abril de 2025}
\release{1}
\author{Alejandro}
\newcommand{\sphinxlogo}{\vbox{}}
\renewcommand{\releasename}{Versión}
\makeindex
\begin{document}

\ifdefined\shorthandoff
  \ifnum\catcode`\=\string=\active\shorthandoff{=}\fi
  \ifnum\catcode`\"=\active\shorthandoff{"}\fi
\fi

\pagestyle{empty}
\sphinxmaketitle
\pagestyle{plain}
\sphinxtableofcontents
\pagestyle{normal}
\phantomsection\label{\detokenize{index::doc}}


\sphinxstepscope


\chapter{proyect}
\label{\detokenize{modules:proyect}}\label{\detokenize{modules::doc}}
\sphinxstepscope


\section{proyect package}
\label{\detokenize{proyect:proyect-package}}\label{\detokenize{proyect::doc}}

\subsection{Submodules}
\label{\detokenize{proyect:submodules}}

\subsection{proyect.APP module}
\label{\detokenize{proyect:proyect-app-module}}

\subsection{proyect.Clientes module}
\label{\detokenize{proyect:module-proyect.Clientes}}\label{\detokenize{proyect:proyect-clientes-module}}\index{module@\spxentry{module}!proyect.Clientes@\spxentry{proyect.Clientes}}\index{proyect.Clientes@\spxentry{proyect.Clientes}!module@\spxentry{module}}\index{Cliente (clase en proyect.Clientes)@\spxentry{Cliente}\spxextra{clase en proyect.Clientes}}

\begin{fulllineitems}
\phantomsection\label{\detokenize{proyect:proyect.Clientes.Cliente}}
\pysigstartsignatures
\pysiglinewithargsret{\sphinxbfcode{\sphinxupquote{class\DUrole{w}{ }}}\sphinxcode{\sphinxupquote{proyect.Clientes.}}\sphinxbfcode{\sphinxupquote{Cliente}}}{\sphinxparam{\DUrole{n}{id\_cliente}\DUrole{p}{:}\DUrole{w}{ }\DUrole{n}{str}}\sphinxparamcomma \sphinxparam{\DUrole{n}{nombre}\DUrole{p}{:}\DUrole{w}{ }\DUrole{n}{str}}\sphinxparamcomma \sphinxparam{\DUrole{n}{email}\DUrole{p}{:}\DUrole{w}{ }\DUrole{n}{str}}\sphinxparamcomma \sphinxparam{\DUrole{n}{direccion}\DUrole{p}{:}\DUrole{w}{ }\DUrole{n}{str}}}{}
\pysigstopsignatures
\sphinxAtStartPar
Bases: \sphinxcode{\sphinxupquote{object}}
\index{get\_direccion() (método de proyect.Clientes.Cliente)@\spxentry{get\_direccion()}\spxextra{método de proyect.Clientes.Cliente}}

\begin{fulllineitems}
\phantomsection\label{\detokenize{proyect:proyect.Clientes.Cliente.get_direccion}}
\pysigstartsignatures
\pysiglinewithargsret{\sphinxbfcode{\sphinxupquote{get\_direccion}}}{}{}
\pysigstopsignatures
\sphinxAtStartPar
Cliente get\_dirrecion : Es para mostrar la dirrecion del cliente de forma correcta , esto es necesario ya que los atributos los tenemos en privados
\begin{description}
\sphinxlineitem{Returns:}
\sphinxAtStartPar
direccion (str): Dirrecion del cliente

\end{description}

\end{fulllineitems}

\index{get\_email() (método de proyect.Clientes.Cliente)@\spxentry{get\_email()}\spxextra{método de proyect.Clientes.Cliente}}

\begin{fulllineitems}
\phantomsection\label{\detokenize{proyect:proyect.Clientes.Cliente.get_email}}
\pysigstartsignatures
\pysiglinewithargsret{\sphinxbfcode{\sphinxupquote{get\_email}}}{}{}
\pysigstopsignatures
\sphinxAtStartPar
Cliente get\_email : Es para mostrar el correo del cliente de forma correcta , esto es necesario ya que los atributos los tenemos en privados
\begin{description}
\sphinxlineitem{Returns:}
\sphinxAtStartPar
email (str) : Dirrecion de correo del cliente

\end{description}

\end{fulllineitems}

\index{get\_id\_cliente() (método de proyect.Clientes.Cliente)@\spxentry{get\_id\_cliente()}\spxextra{método de proyect.Clientes.Cliente}}

\begin{fulllineitems}
\phantomsection\label{\detokenize{proyect:proyect.Clientes.Cliente.get_id_cliente}}
\pysigstartsignatures
\pysiglinewithargsret{\sphinxbfcode{\sphinxupquote{get\_id\_cliente}}}{}{}
\pysigstopsignatures
\sphinxAtStartPar
Cliente get\_id\_cliente : Es para mostrar el id del cliente de forma correcta , esto es necesario ya que los atributos los tenemos en privados
\begin{description}
\sphinxlineitem{Returns:}
\sphinxAtStartPar
id\_cliente (str) : Es el indentificador de cada cliente ; es unico para cada cliente

\end{description}

\end{fulllineitems}

\index{get\_nombre() (método de proyect.Clientes.Cliente)@\spxentry{get\_nombre()}\spxextra{método de proyect.Clientes.Cliente}}

\begin{fulllineitems}
\phantomsection\label{\detokenize{proyect:proyect.Clientes.Cliente.get_nombre}}
\pysigstartsignatures
\pysiglinewithargsret{\sphinxbfcode{\sphinxupquote{get\_nombre}}}{}{}
\pysigstopsignatures
\sphinxAtStartPar
Cliente get\_nombre : Es para mostrar el nombre del cliente de forma correcta , esto es necesario ya que los atributos los tenemos en privados
\begin{description}
\sphinxlineitem{Returns:}
\sphinxAtStartPar
nombre (str) : Nombre del cliente puede ser principalmente sin apellidos pero en caso necesario se pueden meter

\end{description}

\end{fulllineitems}

\index{set\_direccion() (método de proyect.Clientes.Cliente)@\spxentry{set\_direccion()}\spxextra{método de proyect.Clientes.Cliente}}

\begin{fulllineitems}
\phantomsection\label{\detokenize{proyect:proyect.Clientes.Cliente.set_direccion}}
\pysigstartsignatures
\pysiglinewithargsret{\sphinxbfcode{\sphinxupquote{set\_direccion}}}{\sphinxparam{\DUrole{n}{direccion}}}{}
\pysigstopsignatures
\sphinxAtStartPar
Cliente set\_dirrecion : Es para cambiar la dirrecion del cliente para esto se le pasa la nueva dirrecion y la modifica , esto es necesario ya que los atributos los tenemos en privados

\end{fulllineitems}

\index{set\_email() (método de proyect.Clientes.Cliente)@\spxentry{set\_email()}\spxextra{método de proyect.Clientes.Cliente}}

\begin{fulllineitems}
\phantomsection\label{\detokenize{proyect:proyect.Clientes.Cliente.set_email}}
\pysigstartsignatures
\pysiglinewithargsret{\sphinxbfcode{\sphinxupquote{set\_email}}}{\sphinxparam{\DUrole{n}{email}}}{}
\pysigstopsignatures
\sphinxAtStartPar
Cliente set\_email : Es para cambiar el correo del cliente para esto se le pasa el nuevo correo y lo modifica , esto es necesario ya que los atributos los tenemos en privados

\end{fulllineitems}

\index{set\_id\_cliente() (método de proyect.Clientes.Cliente)@\spxentry{set\_id\_cliente()}\spxextra{método de proyect.Clientes.Cliente}}

\begin{fulllineitems}
\phantomsection\label{\detokenize{proyect:proyect.Clientes.Cliente.set_id_cliente}}
\pysigstartsignatures
\pysiglinewithargsret{\sphinxbfcode{\sphinxupquote{set\_id\_cliente}}}{\sphinxparam{\DUrole{n}{id\_cliente}}}{}
\pysigstopsignatures
\sphinxAtStartPar
Cliente set\_id\_cliente : Es para cambiar el id del cliente para esto se le pasa el nuevo id y lo modifica , esto es necesario ya que los atributos los tenemos en privados

\end{fulllineitems}

\index{set\_nombre() (método de proyect.Clientes.Cliente)@\spxentry{set\_nombre()}\spxextra{método de proyect.Clientes.Cliente}}

\begin{fulllineitems}
\phantomsection\label{\detokenize{proyect:proyect.Clientes.Cliente.set_nombre}}
\pysigstartsignatures
\pysiglinewithargsret{\sphinxbfcode{\sphinxupquote{set\_nombre}}}{\sphinxparam{\DUrole{n}{nombre}}}{}
\pysigstopsignatures
\sphinxAtStartPar
Cliente set\_nombre : Es para cambiar el nombre del cliente para esto se le pasa el nuevo nombre y lo modifica , esto es necesario ya que los atributos los tenemos en privados

\end{fulllineitems}


\end{fulllineitems}



\subsection{proyect.Pedido module}
\label{\detokenize{proyect:proyect-pedido-module}}

\subsection{proyect.Productos module}
\label{\detokenize{proyect:module-proyect.Productos}}\label{\detokenize{proyect:proyect-productos-module}}\index{module@\spxentry{module}!proyect.Productos@\spxentry{proyect.Productos}}\index{proyect.Productos@\spxentry{proyect.Productos}!module@\spxentry{module}}\index{Productos (clase en proyect.Productos)@\spxentry{Productos}\spxextra{clase en proyect.Productos}}

\begin{fulllineitems}
\phantomsection\label{\detokenize{proyect:proyect.Productos.Productos}}
\pysigstartsignatures
\pysiglinewithargsret{\sphinxbfcode{\sphinxupquote{class\DUrole{w}{ }}}\sphinxcode{\sphinxupquote{proyect.Productos.}}\sphinxbfcode{\sphinxupquote{Productos}}}{\sphinxparam{\DUrole{n}{id\_producto}\DUrole{p}{:}\DUrole{w}{ }\DUrole{n}{str}}\sphinxparamcomma \sphinxparam{\DUrole{n}{nombre}\DUrole{p}{:}\DUrole{w}{ }\DUrole{n}{str}}\sphinxparamcomma \sphinxparam{\DUrole{n}{precio}\DUrole{p}{:}\DUrole{w}{ }\DUrole{n}{float}}\sphinxparamcomma \sphinxparam{\DUrole{n}{stock}\DUrole{p}{:}\DUrole{w}{ }\DUrole{n}{int}}}{}
\pysigstopsignatures
\sphinxAtStartPar
Bases: \sphinxcode{\sphinxupquote{object}}
\index{get\_id\_productos() (método de proyect.Productos.Productos)@\spxentry{get\_id\_productos()}\spxextra{método de proyect.Productos.Productos}}

\begin{fulllineitems}
\phantomsection\label{\detokenize{proyect:proyect.Productos.Productos.get_id_productos}}
\pysigstartsignatures
\pysiglinewithargsret{\sphinxbfcode{\sphinxupquote{get\_id\_productos}}}{}{}
\pysigstopsignatures\begin{description}
\sphinxlineitem{Returns:}
\sphinxAtStartPar
str: devueve el id producto

\end{description}

\end{fulllineitems}

\index{get\_nombre() (método de proyect.Productos.Productos)@\spxentry{get\_nombre()}\spxextra{método de proyect.Productos.Productos}}

\begin{fulllineitems}
\phantomsection\label{\detokenize{proyect:proyect.Productos.Productos.get_nombre}}
\pysigstartsignatures
\pysiglinewithargsret{\sphinxbfcode{\sphinxupquote{get\_nombre}}}{}{}
\pysigstopsignatures\begin{description}
\sphinxlineitem{Returns:}
\sphinxAtStartPar
str: devuelve el nombre del producto

\end{description}

\end{fulllineitems}

\index{get\_precio() (método de proyect.Productos.Productos)@\spxentry{get\_precio()}\spxextra{método de proyect.Productos.Productos}}

\begin{fulllineitems}
\phantomsection\label{\detokenize{proyect:proyect.Productos.Productos.get_precio}}
\pysigstartsignatures
\pysiglinewithargsret{\sphinxbfcode{\sphinxupquote{get\_precio}}}{}{}
\pysigstopsignatures\begin{description}
\sphinxlineitem{Returns:}
\sphinxAtStartPar
float: devuelve el precio del producto

\end{description}

\end{fulllineitems}

\index{get\_stock() (método de proyect.Productos.Productos)@\spxentry{get\_stock()}\spxextra{método de proyect.Productos.Productos}}

\begin{fulllineitems}
\phantomsection\label{\detokenize{proyect:proyect.Productos.Productos.get_stock}}
\pysigstartsignatures
\pysiglinewithargsret{\sphinxbfcode{\sphinxupquote{get\_stock}}}{}{}
\pysigstopsignatures\begin{description}
\sphinxlineitem{Returns:}
\sphinxAtStartPar
int: devuelve el stock del producto

\end{description}

\end{fulllineitems}

\index{set\_id\_productos() (método de proyect.Productos.Productos)@\spxentry{set\_id\_productos()}\spxextra{método de proyect.Productos.Productos}}

\begin{fulllineitems}
\phantomsection\label{\detokenize{proyect:proyect.Productos.Productos.set_id_productos}}
\pysigstartsignatures
\pysiglinewithargsret{\sphinxbfcode{\sphinxupquote{set\_id\_productos}}}{\sphinxparam{\DUrole{n}{id\_producto}}}{}
\pysigstopsignatures\begin{description}
\sphinxlineitem{Args:}
\sphinxAtStartPar
id\_producto (str): es el id del producto que se le pasa para que modifique el parametro

\end{description}

\end{fulllineitems}

\index{set\_nombre() (método de proyect.Productos.Productos)@\spxentry{set\_nombre()}\spxextra{método de proyect.Productos.Productos}}

\begin{fulllineitems}
\phantomsection\label{\detokenize{proyect:proyect.Productos.Productos.set_nombre}}
\pysigstartsignatures
\pysiglinewithargsret{\sphinxbfcode{\sphinxupquote{set\_nombre}}}{\sphinxparam{\DUrole{n}{nombre}}}{}
\pysigstopsignatures\begin{description}
\sphinxlineitem{Args:}
\sphinxAtStartPar
nombre (str): es el nombre del producto que se le pasa para que modifique el parametro

\end{description}

\end{fulllineitems}

\index{set\_precio() (método de proyect.Productos.Productos)@\spxentry{set\_precio()}\spxextra{método de proyect.Productos.Productos}}

\begin{fulllineitems}
\phantomsection\label{\detokenize{proyect:proyect.Productos.Productos.set_precio}}
\pysigstartsignatures
\pysiglinewithargsret{\sphinxbfcode{\sphinxupquote{set\_precio}}}{\sphinxparam{\DUrole{n}{precio}}}{}
\pysigstopsignatures\begin{description}
\sphinxlineitem{Args:}
\sphinxAtStartPar
precio (float): es el precio del producto que se le pasa para cambiarlo

\end{description}

\end{fulllineitems}

\index{set\_stock() (método de proyect.Productos.Productos)@\spxentry{set\_stock()}\spxextra{método de proyect.Productos.Productos}}

\begin{fulllineitems}
\phantomsection\label{\detokenize{proyect:proyect.Productos.Productos.set_stock}}
\pysigstartsignatures
\pysiglinewithargsret{\sphinxbfcode{\sphinxupquote{set\_stock}}}{\sphinxparam{\DUrole{n}{stock}}}{}
\pysigstopsignatures\begin{description}
\sphinxlineitem{Args:}
\sphinxAtStartPar
stock (int): es el stock del producto que se le pasa para que modifique el parametro

\end{description}

\end{fulllineitems}


\end{fulllineitems}



\subsection{proyect.Productos\_Digitales module}
\label{\detokenize{proyect:proyect-productos-digitales-module}}

\subsection{proyect.reseña module}
\label{\detokenize{proyect:proyect-resena-module}}

\subsection{Module contents}
\label{\detokenize{proyect:module-proyect}}\label{\detokenize{proyect:module-contents}}\index{module@\spxentry{module}!proyect@\spxentry{proyect}}\index{proyect@\spxentry{proyect}!module@\spxentry{module}}

\chapter{Indices and tables}
\label{\detokenize{index:indices-and-tables}}\begin{itemize}
\item {} 
\sphinxAtStartPar
\DUrole{xref,std,std-ref}{genindex}

\item {} 
\sphinxAtStartPar
\DUrole{xref,std,std-ref}{modindex}

\item {} 
\sphinxAtStartPar
\DUrole{xref,std,std-ref}{search}

\end{itemize}


\renewcommand{\indexname}{Índice de Módulos Python}
\begin{sphinxtheindex}
\let\bigletter\sphinxstyleindexlettergroup
\bigletter{p}
\item\relax\sphinxstyleindexentry{proyect}\sphinxstyleindexpageref{proyect:\detokenize{module-proyect}}
\item\relax\sphinxstyleindexentry{proyect.Clientes}\sphinxstyleindexpageref{proyect:\detokenize{module-proyect.Clientes}}
\item\relax\sphinxstyleindexentry{proyect.Productos}\sphinxstyleindexpageref{proyect:\detokenize{module-proyect.Productos}}
\end{sphinxtheindex}

\renewcommand{\indexname}{Índice}
\printindex
\end{document}